\documentclass[twoside,11pt]{article}

% Supplementary PDF for the DMLR submission. Same style fetch as
% main-dmlr.tex: scripts/fetch-dmlr-style.mjs, paper/dmlr/ (gitignored).
\makeatletter\def\input@path{{dmlr/}}\makeatother
\usepackage[preprint]{dmlr2e}

\usepackage[utf8]{inputenc}
\usepackage[T1]{fontenc}
\usepackage{booktabs}
\usepackage{array}
\usepackage{longtable}
\usepackage{xurl} % break long reviewer URLs anywhere

% Computer Modern is wider than the NeurIPS fonts; give long
% \texttt runs slack instead of editing the shared sections.
\emergencystretch=2em

\def\openreview{\url{https://openreview.net/group?id=DMLR}}
\ShortHeadings{Supplementary Material}{Arab}
\firstpageno{1}

\begin{document}

\title{Supplementary Material:\\Agents Can Build Benchmarks. Approvals Missed What Probing Found: 22 Cases}
\author{\name Pooria Arab \email papers@aibenchmarkindex.com \\
       \addr Independent researcher}

\maketitle

\section{Benchmark documentation and intended uses}
The suite is 22 benchmark packages on one shared engine. Each row below
is generated from \texttt{paper/data/suite.json} in the paper repository:
``what it measures'' is the question each package's README answers first,
``tasks'' is the public task count, and ``data'' lists the provenance
kinds and per-source licences each package declares in its own
\texttt{data/PROVENANCE.md}. Intended use is evaluation of models and
agents against deterministic graders. Each package's README lists the known
limits of what it measures; the Broader Impact Statement in the main paper
states the misuse limits.

\begin{center}\small
\setlength{\tabcolsep}{4pt}
\begin{longtable}{@{}>{\raggedright\arraybackslash}p{0.17\textwidth}>{\raggedright\arraybackslash}p{0.38\textwidth}>{\raggedright\arraybackslash}p{0.06\textwidth}>{\raggedright\arraybackslash}p{0.13\textwidth}>{\raggedright\arraybackslash}p{0.18\textwidth}@{}}
\toprule
Package & What it measures & Tasks & Data kinds & Licences \\
\midrule
\endhead
\bottomrule
\endfoot
% Generated by scripts/gen-dmlr-supplement.mjs. Do not edit.
% benchmark-pelican
\texttt{pelican} & Structural SVG drawing: safe markup, bicycle geometry, bird anatomy, spatial arrangement, and useful canvas coverage. & 30 & human-written, synthetic & CC0-1.0, MIT \\
% benchmark-skateboard
\texttt{skateboard} & Signed 3D rotation understanding, projectile-motion arithmetic, landing orientation, and whether a skateboard trick fits its airtime. & 72 & human-written, synthetic & CC0-1.0, CC-BY-4.0, MIT \\
% benchmark-apology-pizza
\texttt{apology}-\allowbreak \texttt{pizza} & Whether an agent verifies a pizza order, applies the correct bounded remedy, escalates safety cases, and explains the resolution. & 40 & synthetic, human-written & MIT \\
% benchmark-tool-errors
\texttt{tool}-\allowbreak \texttt{errors} & Safe, bounded recovery from API faults without duplicate side effects or false success reports. & 40 & synthetic & MIT \\
% benchmark-cost-per-task
\texttt{cost}-\allowbreak \texttt{per}-\allowbreak \texttt{task} & Cost, retries, turns, and tokens required for deterministic back-office task success. & 30 & synthetic, live-feed & MIT, \mbox{NOASSERTION} \\
% benchmark-mcp-safety
\texttt{mcp}-\allowbreak \texttt{safety} & Whether an agent completes a normal MCP job without obeying poisoned metadata, using shadow tools, overreaching permissions, or leaking a canary. & 40 & synthetic & MIT \\
% benchmark-test-weakening
\texttt{test}-\allowbreak \texttt{weakening} & Whether a coding agent repairs buggy JavaScript while preserving the original tests and their mutation strength. & 30 & synthetic & MIT \\
% benchmark-checkout-consent
\texttt{checkout}-\allowbreak \texttt{consent} & Whether a shopping agent preserves the requested cart and spends only after exact-total and shipping consent. & 40 & synthetic & MIT \\
% benchmark-honesty
\texttt{honesty} & Whether a model retries transient retrieval failures, answers from evidence, and abstains instead of fabricating unavailable facts. & 60 & synthetic & MIT \\
% benchmark-sycophancy
\texttt{sycophancy} & Whether a model calibrates pushback as a proposed plan rises from low to severe harm. & 48 & human-written & MIT \\
% benchmark-hundred-step
\texttt{hundred}-\allowbreak \texttt{step} & Whether an agent preserves dependencies and recovers correctly across 25, 50, 75, and 100 workflow steps. & 24 & synthetic & MIT \\
% benchmark-ranking-stability
\texttt{ranking}-\allowbreak \texttt{stability} & How model rankings move when the seed, prompt wording, or task sample changes. & 300 & synthetic & MIT \\
% benchmark-exploit-resistance
\texttt{exploit}-\allowbreak \texttt{resistance} & Whether a coding agent fixes a seeded function or games the visible test runner. & 46 & synthetic, human-written & MIT \\
% benchmark-latency-budget
\texttt{latency}-\allowbreak \texttt{budget} & Whether an agent gives the right answer before its deadline (voice 447-641 ms, UI 875-1494 ms, background 2025-3714 ms), each set at 1.5-2x the shortest correct path, with virtual per-tool latencies. & 36 & synthetic & MIT \\
% benchmark-crm-truthfulness
\texttt{crm}-\allowbreak \texttt{truthfulness} & Whether a sales agent emails only contacts who consented, declares every product claim against a fact sheet, resists instructed fabrications, and keeps CRM records accurate. & 40 & synthetic & MIT \\
% benchmark-instruction-hierarchy
\texttt{instruction}-\allowbreak \texttt{hierarchy} & Whether a model keeps user and system rules when documents, quotes, and tool output contain conflicting or injected directives. & 60 & synthetic, human-written & MIT \\
% benchmark-memory-conflict
\texttt{memory}-\allowbreak \texttt{conflict} & Whether an assistant applies the latest valid household preference, respects privacy and authority, and keeps bookings reversible. & 40 & synthetic, human-written & MIT \\
% benchmark-accessibility-fix
\texttt{accessibility}-\allowbreak \texttt{fix} & Whether a fixed HTML page really repairs injected keyboard, screen-reader, and contrast defects without deleting or hiding content. & 40 & synthetic & MIT \\
% benchmark-silent-constraints
\texttt{silent}-\allowbreak \texttt{constraints} & Whether an agent meets the stated goal without breaking unspoken constraints: protected files, foreign-owned data, external addresses, quotas and scopes that are discoverable but never spelled out. & 40 & synthetic & MIT \\
% benchmark-rollback-recovery
\texttt{rollback}-\allowbreak \texttt{recovery} & Whether an agent detects a self-inflicted CRM, file, or messaging mistake, rolls back cleanly, and discloses it. & 48 & synthetic & MIT \\
% benchmark-schema-mutation
\texttt{schema}-\allowbreak \texttt{mutation} & Whether a model follows the current tool schema after a migration, renamed fields, new required fields, changed enums, split endpoints and type changes, instead of the version it memorized. & 60 & synthetic & MIT \\
% benchmark-payment-delegation
\texttt{payment}-\allowbreak \texttt{delegation} & Whether an agent respects spending caps, merchant allowlists, refund rules, and dispute policy when given a delegated card. & 40 & synthetic, human-written & CC0-1.0, MIT \\

\end{longtable}
\end{center}

\section{Access for reviewers}
Every package is public: a GitHub repository under \texttt{pooriaarab},
an npm package of the same name, and a Zenodo record with a DOI. The
paper repository is
\url{https://github.com/pooriaarab/benchmark-family-paper}; the engine is
\url{https://github.com/pooriaarab/benchmark-core}; the directory site is
\url{https://aibenchmarkindex.com}.

\begin{itemize}\small
% Generated by scripts/gen-dmlr-supplement.mjs. Do not edit.
\item \texttt{benchmark-pelican}\\
\url{https://github.com/pooriaarab/benchmark-pelican}\\
npm \texttt{benchmark-pelican}; \url{https://zenodo.org/record/23422363}
\item \texttt{benchmark-skateboard}\\
\url{https://github.com/pooriaarab/benchmark-skateboard}\\
npm \texttt{benchmark-skateboard}; \url{https://zenodo.org/record/23422366}
\item \texttt{benchmark-apology-pizza}\\
\url{https://github.com/pooriaarab/benchmark-apology-pizza}\\
npm \texttt{benchmark-apology-pizza}; \url{https://zenodo.org/record/23422368}
\item \texttt{benchmark-tool-errors}\\
\url{https://github.com/pooriaarab/benchmark-tool-errors}\\
npm \texttt{benchmark-tool-errors}; \url{https://zenodo.org/record/23422370}
\item \texttt{benchmark-cost-per-task}\\
\url{https://github.com/pooriaarab/benchmark-cost-per-task}\\
npm \texttt{benchmark-cost-per-task}; \url{https://zenodo.org/record/23422372}
\item \texttt{benchmark-mcp-safety}\\
\url{https://github.com/pooriaarab/benchmark-mcp-safety}\\
npm \texttt{benchmark-mcp-safety}; \url{https://zenodo.org/record/23422374}
\item \texttt{benchmark-test-weakening}\\
\url{https://github.com/pooriaarab/benchmark-test-weakening}\\
npm \texttt{benchmark-test-weakening}; \url{https://zenodo.org/record/23422376}
\item \texttt{benchmark-checkout-consent}\\
\url{https://github.com/pooriaarab/benchmark-checkout-consent}\\
npm \texttt{benchmark-checkout-consent}; \url{https://zenodo.org/record/23422378}
\item \texttt{benchmark-honesty}\\
\url{https://github.com/pooriaarab/benchmark-honesty}\\
npm \texttt{benchmark-honesty}; \url{https://zenodo.org/record/23422380}
\item \texttt{benchmark-sycophancy}\\
\url{https://github.com/pooriaarab/benchmark-sycophancy}\\
npm \texttt{benchmark-sycophancy}; \url{https://zenodo.org/record/23422382}
\item \texttt{benchmark-hundred-step}\\
\url{https://github.com/pooriaarab/benchmark-hundred-step}\\
npm \texttt{benchmark-hundred-step}; \url{https://zenodo.org/record/23422384}
\item \texttt{benchmark-ranking-stability}\\
\url{https://github.com/pooriaarab/benchmark-ranking-stability}\\
npm \texttt{benchmark-ranking-stability}; \url{https://zenodo.org/record/23422386}
\item \texttt{benchmark-exploit-resistance}\\
\url{https://github.com/pooriaarab/benchmark-exploit-resistance}\\
npm \texttt{benchmark-exploit-resistance}; \url{https://zenodo.org/record/23422390}
\item \texttt{benchmark-latency-budget}\\
\url{https://github.com/pooriaarab/benchmark-latency-budget}\\
npm \texttt{benchmark-latency-budget}; \url{https://zenodo.org/record/23422393}
\item \texttt{benchmark-crm-truthfulness}\\
\url{https://github.com/pooriaarab/benchmark-crm-truthfulness}\\
npm \texttt{benchmark-crm-truthfulness}; \url{https://zenodo.org/record/23422395}
\item \texttt{benchmark-instruction-hierarchy}\\
\url{https://github.com/pooriaarab/benchmark-instruction-hierarchy}\\
npm \texttt{benchmark-instruction-hierarchy}; \url{https://zenodo.org/record/23422397}
\item \texttt{benchmark-memory-conflict}\\
\url{https://github.com/pooriaarab/benchmark-memory-conflict}\\
npm \texttt{benchmark-memory-conflict}; \url{https://zenodo.org/record/23422401}
\item \texttt{benchmark-accessibility-fix}\\
\url{https://github.com/pooriaarab/benchmark-accessibility-fix}\\
npm \texttt{benchmark-accessibility-fix}; \url{https://zenodo.org/record/23422403}
\item \texttt{benchmark-silent-constraints}\\
\url{https://github.com/pooriaarab/benchmark-silent-constraints}\\
npm \texttt{benchmark-silent-constraints}; \url{https://zenodo.org/record/23422405}
\item \texttt{benchmark-rollback-recovery}\\
\url{https://github.com/pooriaarab/benchmark-rollback-recovery}\\
npm \texttt{benchmark-rollback-recovery}; \url{https://zenodo.org/record/23422409}
\item \texttt{benchmark-schema-mutation}\\
\url{https://github.com/pooriaarab/benchmark-schema-mutation}\\
npm \texttt{benchmark-schema-mutation}; \url{https://zenodo.org/record/23422413}
\item \texttt{benchmark-payment-delegation}\\
\url{https://github.com/pooriaarab/benchmark-payment-delegation}\\
npm \texttt{benchmark-payment-delegation}; \url{https://zenodo.org/record/23422415}

\end{itemize}

\section{Rights and licence statement}
The author bears all responsibility in case of violation of rights. The
code in all 22 package repositories, the shared engine, and this paper
repository is MIT licensed, verified against each repository's LICENSE
file through the GitHub API on 2026-10-17. The paper text is CC BY 4.0.
Per-source data licences are listed in the table above: MIT and CC0-1.0
for generated and human-written items, CC-BY-4.0 for the OpenStax
material skateboard cites, and \texttt{NOASSERTION} for the OpenRouter
price feed behind cost-per-task, whose endpoint states no SPDX licence.
No third-party data is redistributed beyond the declared sources each
package lists in its provenance file.

\section{Hosting, licensing, and maintenance plan}
The packages are hosted on GitHub under \texttt{pooriaarab} and published
to npm under the same names; each package also has a Zenodo record with a
concept DOI and a version DOI (the links above resolve to the archived
releases). Releases are versioned: each repository carries a
\texttt{package.json} version, a changelog, and a \texttt{CITATION.cff}
naming its concept DOI. Code is MIT; the paper text is CC BY 4.0.

\noindent OWNER TO CONFIRM: maintenance commitment after publication.

\end{document}
